\section{Exact arithmetic and the stored state}\label{app:arithmetic}

The budget $E_k\le n$ bounds the number of digit assignments and their
length. We verify that it also bounds their bit cost, and that the
adversary's data structure remains of polynomial size during the first
$Q(n)$ comparisons.
Computation is measured on a fixed deterministic multitape Turing
machine, with integers represented in binary.

Write $S_k$ for the number of digit choices, $t_k$ for the number of
sweeps, and $b_k$ for the bound on retained vertices in
\eqref{eq:local-size}, with the parameters chosen in
Appendix~\ref{app:certificates}. The enumeration bound is
\begin{equation}\label{eq:enumeration-budget}
E_k=S_k^{b_k t_k}.
\end{equation}

\subsection{Choosing the parameter}

A fixed program generates the parameters and $E_k$ for
$k=64,128,\ldots$, writing each completed tuple in binary. On input
size $n$, run this program for $n$ steps and choose the largest completed
$k$ with $E_k\le n$. If none qualifies, answer from the fixed label
order. The simulation writes at most $n$ bits, so parsing its output
and selecting a tuple also take polynomial time.

For every fixed $k$, both its generation time and $E_k$ are finite.
That tuple therefore qualifies for all sufficiently large $n$, proving
$k(n)\to\infty$.

\subsection{Precision in the interval calculation}\label{sec:digit-arithmetic}

Suppose the stored endpoints have at most $\Lambda$ fractional binary
digits. An update takes maxima and minima, forms an integer linear
combination, and divides by the power of two $S$. It therefore adds at
most $\log S$ fractional digits. If there are $b'$ interior vertices,
each endpoint calculation makes $b't$ updates. The two calculations share
digits but use their own values, so the precision in either is at most
\begin{equation}\label{eq:digit-precision}
\Lambda+b't\log S\le\Lambda+\log E.
\end{equation}
This proves the precision claim of Theorem~\ref{thm:local}. All computed
positions remain in $[0,1]$, so their integer numerators need at most one
additional bit. Digits and counts have at most $\log S$ and $\log E+1$
bits, respectively.

\subsection{Precision of new anchors}

If a nondegenerate interval has dyadic endpoints with at most $F$
fractional digits, its length is at least $2^{-F}$. The rounding rule
of Lemma~\ref{lem:rounding} therefore uses a grid exponent at most $F$,
and its anchor has at most $F$ fractional digits.

For a certified central interval, \eqref{eq:digit-precision} gives
$F\le\Lambda+\log E_k$. Forming the middle half of a stored interval
gives endpoints with at most $\Lambda+2$ digits. Since $E_k\ge4$, both
choices produce anchors with at most $\Lambda+\log E_k$ digits. This
also covers the direct middle-half choice after a failed comparison
certificate. The update of neighboring bounds and
\texttt{Normalize} take maxima and minima of existing endpoints and the
chosen anchor, adding no precision.

All initial endpoints are $0$ or $1$. After $P$ refinements, induction
therefore gives $\Lambda\le P\log E_k$. By
Lemma~\ref{lem:rare-refinements}, the first $Q(n)$ completed calls to
\texttt{Compare} account for $O(Q(n)/k)$ refinements. An unfinished call
adds at most $n+2$ by Proposition~\ref{prop:legality}. Thus, during these
first $Q(n)$ calls,
\begin{equation}\label{eq:refinement-runtime}
P=O\bigl(Q(n)/k+n\bigr),
\end{equation}
and consequently
\begin{equation}\label{eq:precision}
\Lambda\le P\log E_k=O(n\log^2 n).
\end{equation}
Temporary endpoints in the next calculation add at most another
$\log E_k$ digits.

\subsection{Inventory and work per call}

Each refinement creates at most one node in the adversary's tree. There are thus
at most $P+1$ nodes, $n$ item records, and at most $P$ stored child anchors.
An item belongs to exactly one node graph, so their total numbers of
vertices and arcs are at most $n$ and $n^2$, respectively. The recorded
prefix contains at most $Q(n)$ comparisons. Labels, pointers, depths,
and counters have polynomial length, as do all dyadic endpoints by
\eqref{eq:precision}.

For one certificate, the algorithm enumerates at most $E_k\le n$ digit
assignments. Each assignment uses two endpoint calculations, each with
at most $b_k t_k\le\log E_k\le\log n$ updates. The boundary vertices
remain fixed. Maxima, minima, integer multiplication, and division by
powers of two operate on polynomially many bits. Central certification
sorts at most $n$ lower and upper endpoints and counts the intervals
inside its candidate. Rounding computes a grid exponent and a ceiling
using the same bounded-precision integers. Standard binary arithmetic
therefore gives a fixed polynomial bound for each operation.

Before selecting an answer, a call makes at most two failure refinements
and three calls to \texttt{Select}. Failure refinements use the middle
half directly. After recording the answer, the single structural repair
makes at most $n$ refinements and hence at most $n$ calls to
\texttt{ChooseAnchor}.

The remaining work is polynomial in the stored graph size. Topological
passes normalize intervals and compute directed-path lengths. A finite
graph search finds each retained neighborhood; it stops when either the
prescribed radius or all available vertices have been reached. Routing
uses the adversary's tree, and deterministic choices can be made by scanning
labels. At the horizon, topologically sorting the bounded transcript DAG
produces the fixed ranks for all later answers. Thus the bit-operation
bound has a fixed exponent, independent of $k$ and of the number of
queries already answered.
